\documentclass{article}
\usepackage{tikz}
\usetikzlibrary{angles, quotes}
\usepackage{amssymb}
\usepackage{pgfplots}
\pgfplotsset{width=10cm,compat=1.9}
\usepgfplotslibrary{statistics}
\usepackage[left=2.54cm,right=2.54cm,top=2.54cm,bottom=2.54cm]{geometry}
\usepackage{tasks} %Paquete necesario para la producción de items horizontales para algunos ejercicios de tarea empleando el comando /begin{multicols}}{#}.../end{multicols} para ayudarnos a reducir las páginas
\usepackage{pdflscape} %comando que permite cambiar a orientación horizontal las hojas%
\usepackage{soul} %Paquete para permitir la compilación de acentos usando el comando {\hl{}}}\\
\sethlcolor{yellow(munsell)} %comando para definir el color para el subrayado usando el comando \hl
\usepackage[utf8]{inputenc}
\usepackage[spanish]{babel}
\usepackage{icomma}
\usepackage{ragged2e}
\usepackage{siunitx}
\usepackage{url}
\usepackage[colorlinks=true, urlcolor=blue,  linkcolor=black, citecolor=green]{hyperref}
\usepackage{pdfpages}
\usepackage{blindtext} %Paquete que genera texto ficticio (dummy text) con el comando: \blindtext
\usepackage{cancel} %in the preamble gives you four different modes of striking through: \cancel{text to cancel} draws a diagonal line (slash) through its argument, \bcancel{text to cancel} uses the negative slope (a backslash), \xcancel{text to cancel} draws an X (actually \cancel plus \bcancel), \cancelto{〈value〉}{〈expression〉} draws a diagonal arrow through the 〈expression〉pointing to the 〈value〉 (math-mode only)
\usepackage{csquotes}
\usepackage{afterpage}
\usepackage{parskip} 
\usepackage{float}
\usepackage{enumitem}
\usepackage{multicol}%Paquete que permite la creación aislada de columnas de texto. De esta forma se reduce la cantidad de páginas en nuestro documento
\usepackage{lipsum} %Paquete que genera texto de relleno al igual que \usepackage{blindtext}, se usa con el comando \lipsum[#-#] los #(hashtags) delimitan el número de páginas que deseamos ocupar del paquete lipsum para usarlos como relleno. 

% Definir el comando personalizado para la nota
\newcommand{\mynote}[1]{%
\begin{tikzpicture}[baseline=-0.75ex]
\node[draw, circle, fill=Apple Green!60, inner sep=2pt] (note) {\textbf{Nota:}};
\end{tikzpicture}%
\ \textit{#1}%
}

\newenvironment{Figura}
{\par\medskip\noindent\minipage{\linewidth}}
{\endminipage\par\medskip}
\usepackage{caption}
\usepackage[
backend=biber,
sorting=none,
url=true
]{biblatex} %bibliografía
\addbibresource{biblio.bib}
% Establecer el espacio entre las entradas de la bibliografía
\setlength{\bibitemsep}{1\baselineskip} % Puedes ajustar el valor según tus preferencias
\usepackage{amsmath,amsthm,amssymb,amsfonts}
\usepackage{pifont} %Permite la compilación del comando \xmark: es el símbolo de la tachita
\usepackage{mathtools} %permite la compilación de símbolos para matrices
\usepackage{empheq} %Paquete que se relaciona con \usepackage[most]{tcolorbox} para la creación de cuadros/cajas de colores para delimitar los resultados matemáticos o líneas de texto usando la línea de comando: \begin{empheq}[box={\mymath[colback=orange(webcolor)!70,drop lifted shadow]}]{equation*} \end{empheq}
\usepackage[most]{tcolorbox} %Paquete que se relaciona con \usepackage{empheq} para la creación de cuadros/cajas de colores para delimitar los resultados matemáticos o líneas de texto
\renewcommand{\qedsymbol}{$\blacksquare$}
\usetikzlibrary{positioning,decorations.pathreplacing} %Paquete que permite la compilación de llaves para cuadros sinópticos (brace diagram) 
\usepackage{schemata} %The schemata package is designed just for these brace diagrams
\usepackage{booktabs}

\newcommand\AB[2]{\schema{\schemabox{#1}}{\schemabox{#2}}} %comando o paquete necesario para crear cuadros sinópticos

\newtcbox{\mymath}[1][]{%
nobeforeafter, math upper, tcbox raise base,
enhanced, colframe=blue!30!black,
colback=blue!30, boxrule=1pt,
#1} %Comando importante para encasillar los resultados matemáticos en cajas de diferentes colores, se relaciona con los paquetes \usepackage{empheq} y \usepackage[most]{tcolorbox} 

\newenvironment{sysmatrix}[1]
{\left(\begin{array}{@{}#1@{}}}
{\end{array}\right)}
\newcommand{\ro}[1]{%
\xrightarrow{\mathmakebox[\rowidth]{#1}}%
}
\newlength{\rowidth}% row operation width
\AtBeginDocument{\setlength{\rowidth}{3em}}  %Comando importante para laproducción de lineas de operaciones entre matrices, método de Gauss o Gauss Jordan se relaciona con \newenvironment{sysmatrix}

%formato para cambiar el horario a español
\usepackage[spanish]{datetime2}
\DTMsetdatestyle{spanish}

\renewcommand{\today}{\DTMdisplaydate{\the\year}{\the\month}{\the\day }{-1}}
%%%%%%%%%%%%%%%%%%%%%%%%%%%%%%%%%%%%%%%%%

\usepackage{datetime}
\newcommand{\mycurrenttime}{\xxivtime}

%%%%%%%%%%%%%%%%%%%%%%%%%%%%%%%%%%%%%%%%%%%%%%%%%%%%%%%%%%%%%%%%%%%%%%%%%%%%%%%%%%%%%%%%%%%%%%%%%%%%%%%%%%%%%%%%%%%%%%%%%%%%%%%%%%%%%%%%%%%%%%%%%


%%%%%%%%%%%%%Caja de color para el título%%%%%%%%%%%%%%%%%%%%%%%%%%%%%%%%%%%%%%%%%%%%%%%%%%%%%%

\definecolor{myframecolor}{RGB}{1, 45, 75} % Prussian Blue
\definecolor{myboxcolor}{RGB}{0, 107, 92} % Apple Green

\newtcolorbox{mybox}{
enhanced,
colback=myboxcolor!7, % Color de fondo del cuadro
colframe=myframecolor, % Color del marco del cuadro
arc=0pt,
boxrule=1pt,
borderline west={2mm}{-10mm}{myframecolor}, % Borde en el lado izquierdo
sharp corners=southwest,
width=\linewidth,
before=\par\vspace{\bigskipamount}, % Espacio antes del cuadro
after=\par\vspace{\bigskipamount} % Espacio después del cuadro
}


\newcommand{\euler}{e} %Comando para producir letra e de euler
\newenvironment{remark}{\par\vfill\footnotesize % Vertical white space above the remark and smaller font size
\begin{list}{}{
	\leftmargin=80pt % Indentation on the left
	\rightmargin=60pt}\item\ignorespaces % Indentation on the right
\makebox[-2.5pt]{\begin{tikzpicture}[overlay]
		\node[draw=Horizon!60,line width=2.5pt,circle,fill=Horizon!25,font=\sffamily\bfseries,inner sep=4pt,outer sep=0pt] at (-19pt,5pt){\textcolor{Horizon}{Nota}};\end{tikzpicture}} % Blue Nota in a circle
\advance\baselineskip -1pt}{\end{list}\vskip5pt} % Tighter line spacing and white space after remark
\usepackage{graphicx}

\usepackage{titling}

\input{colores.tex}

%%Producir signos de cita textual``Asignatura''

\newcommand{\xmark}{\ding{55}} %%%comando que genera la tachita

\newenvironment{MyColorPar}[1]{%
\leavevmode\color{#1}\ignorespaces%
}{%
}%

%%%%%%%%%%%%%%%%%%% Título de la tarea, Nombre de alumno

\title{Clase \# 44 - Mecánica Teórica}
\author{\emph{JA} $\boldsymbol{\mid}$}
\date{\today}

%%%%%%%%%%%%%%%%%%% Datos de la Materia

\usepackage{fancyhdr}
\fancypagestyle{plain}{%  the preset of fancyhdr 
\fancyhf{} % clear all header and footer fields
\fancyfoot[R]{
\includegraphics[width=2cm]{LogoFCFMBUAP (1).png}}
\fancyfoot[L]{{\bfseries{\thedate{}}} a las {\bfseries{\mycurrenttime{}}} horas{} (GMT-6, H. Puebla de Zaragoza, Pue)}
\fancyhead[L]{Mecánica teórica}
\fancyhead[R]{\theauthor}
}
\makeatletter
\def\@maketitle{%
\newpage
\null
\vskip 1em%
\begin{center}%
\let \footnote \thanks
{\LARGE \@title \par}%
\vskip 1em%
%{\large \@date}%
\end{center}%
\par
\vskip 1em}
\makeatother

\usepackage{lipsum}  
\usepackage{cmbright}


%%%%%%%%%%%%%%%%%%%%%%%%%%%%%%%%%%%%%%%%%%%%%%%%%%%%%%%%%%%%%%%%%%%%%%%%%%%%%%%%%%%%%%%%%%%%%%%%%%%%%%%%%%%%%%%%%%%%%%%%%%%%%%%%%%%%%%%%%%%%%%%%%%%%%%%%%%%%%%%%%%%%%%%%%%%%%%%%%%%%%%%%%%%%%%%%%%%%%%

\begin{document}



\maketitle


%%%%%%%%%%%%%%%%%%% Datos del profesor, campus y aula

\begin{mybox}
\noindent\begin{tabular}{@{}ll}
	{\bfseries{Profesor}} & Dr. Gilberto Silva Ortigoza\\
	\textcolor{prussianblue}{\bfseries{Campus}} / \textcolor{trueblue}{\bfseries{Facultad}}  & \textcolor{prussianblue}{\bfseries{C.U. BUAP}} /  \textcolor{trueblue}{\bfseries{FCFM}} \\
	\textcolor{prussianblue}{\bfseries{Edificio}} /  \textcolor{trueblue}{\bfseries{EMA4}}    & 	\textcolor{prussianblue}{\bfseries{FM9}} /  \textcolor{trueblue}{\bfseries{302}}   
\end{tabular} 
\end{mybox} \vspace{1cm}

\tableofcontents

\vspace{2cm}

\section{Clase del día de hoy}


Tenemos nosotros: 


\[ \vec{L} \equiv \vec{r} \times \vec{p} \]


Tenemos:

\[ \left[ L_{i},L_{j}   \right] = \epsilon_{ijk}L_{k}  \]

\begin{center}
\includegraphics[width=1\linewidth]{1.png} % <-- Sustituye esta ruta
\\[2mm]
\textbf{Figure 1}
\end{center}	



Tenemos:


\[ \vec{A} = A_{x}\hat{x} + A_{y}\hat{y}+A_{z}\hat{z} \]


\[ vec{B} = B_{x}\hat{x}+B_{y}\hat{y}+B_{z}\hat{z} \]


Tenemos:


\[ \vec{A} \cdot \vec{B} = A_{x}B_{x}+A_{y}B_{y}+A_{z}B_{z}  \]

Tenemos:


\[  = A_{1}B_{1}+ A_{2}B_{2} +A_{3}B_{3} \]


\[ = \sum\limits_{i=1}^{3}  A_{i}B_{i} = A_{i}B_{i}  \]


Tenemos:


\[  \vec{A} \cdot \vec{B} = A_{i}B_{i}, \quad i=1.2.3 \]


Tenemos:


\begin{equation}
\begin{split}
\vec{A} \times \vec{B}  &= \left| \begin{array}{ccc}
\hat{x} & \hat{y} & \hat{z}\\
A_{x} & A_{y} & A_{z}\\
B_{x} & B_{y} & B_{z}
\end{array}\right| \\
&= \hat{x} (A_{y}B_{z}-A_{z}B_{y}) - \hat{y} (A_{x}B_{z}-A_{z}B_{x}) + \hat{z} (A_{x}B_{y}-A_{y}B_{x})
\end{split}
\end{equation}




Tenemos: 

\begin{center}
\includegraphics[width=1\linewidth]{2.png} % <-- Sustituye esta ruta
\\[2mm]
\textbf{Figure 2}
\end{center}	


Tenemos:


\[ ( \vec{A} \times \vec{B} )_{i} = \epsilon_{ijk} A_{j}B_{k} \]


\[  (\vec{A} \times \vec{B})_{1} = \epsilon_{ijk}A_{j}B_{k}  \]

Tenemos:


\[ \epsilon_{123} = 1, \epsilon_{213} = -\epsilon_{123} \]


Siguiendo estas reglitas podremos desmostrar:


\textbf{Cositas importantes}:

\[ \vec{A} \cdot \vec{B} = A_{i}B_{i}  \]

\[ (\vec{A} \times \vec{B})_{i} = \epsilon_{ijk} A_{j}B_{k}  \]

\[ L_{i} = (\vec{r} \times \vec{p})_{i} = \epsilon_{ijk} \times_{j}p_{k} \]


Entonces

Cada vez que aparecen dos índices repetidos existe una sumatoria.

Tenemos:


\[ L_{i} = (\vec{r} \times \vec{p})_{i} = \epsilon_{is\ell} \times_{s}p_{\ell} \]


\[  L_{j} = (\vec{r}\times \vec{p})_{j} = \epsilon_{jnm} \times_{n}p_{m}  \]


Tenemos:


\[  \left[  L_{i}, L_{j}  \right] = \left[  \epsilon_{is\ell} \times_{s}p_{\ell}, \epsilon_{jnm}\times_{n}p_{m}  \right] \]

\[ = \epsilon_{is\ell}\epsilon_{jnm} \left[ \times_{s} p_{\ell}, \times_{n}p_{m}  \right] \]


Tenemos:



\[ = \epsilon_{is\ell} \epsilon_{jnm} \left\{  \times_{n} \left[ \times_{s}p_{\ell}, p_{m}  \right] + \left[ \times_{s}p_{\ell}, \times_{n}  \right] p_{m}  \right\} \]



\[ \epsilon_{is\ell} \epsilon_{jnm} \left\{   \times_{n}\times_{s}\underbrace{\left[  p_{\ell}, p_{m}  \right]}_{=0} + \times x_{n} \underbrace{\left[ \times_{s},p_{m}  \right]}_{=\delta_{sm}} p_{\ell} + \times_{s} \underbrace{\left[  p_{\ell}, \times_{n} \right]}_{=-\delta_{\ell n}}p_{m} + \underbrace{\left[ x_{s}, x_{n}  \right]}_{=0} p_{\ell}p_{m} \right\}  \]

\textbf{Corchetes fundamentales de Poisson}

\begin{center}
\includegraphics[width=1\linewidth]{3.png} % <-- Sustituye esta ruta
\\[2mm]
\textbf{Figure 3}
\end{center}	

Tenemos:


\[  \epsilon_{is\ell} \epsilon_{jns} = -\epsilon_{si\ell} \epsilon_{sjn} \]


\[  = - \left[  \delta_{ij}\delta_{\ell n} - \delta_{in}\delta_{\ell j}  \right] \]

Si nosotros tenemos lo siguiente (\textbf{primer término}):


\[ \epsilon_{is\ell} \epsilon_{jns} \times_{n}p_{\ell} \]


\[ = - \times_{n} p_{\ell} \left[ \delta_{ij} \delta_{\ell n} - \delta_{in} \delta_{\ell j}    \right] \]

\[  = -\times_{n}p_{n} \delta_{ij}+\times_{i}p_{j} \]


Tenemos:


\[ \epsilon_{is\ell} \epsilon_{j\ell m} \times_{s}p_{m} = -\epsilon_{\ell is} \epsilon_{\ell jm} \times_{s}p_{m}  \]


\[ = -\times_{s}p_{m} \left[   \delta{ij}\delta_{sm}-\delta_{im}\delta_{sj} \right]  \]


\[  = -\times_{m}p_{m}\delta_{ij}+\times_{j}p_{i}\]


Tenemos:


\[  \left[  L_{i}, L_{j}  \right] = - \times_{m}p_{m}\delta_{ij}+\times_{i}p_{j} + \times_{m}p_{m}\delta_{ij}-\times_{j}p_{i} \]

Tenemos:


\[ \vec{A} \times\vec{B} = A_{i}B_{i}=A_{n}B_{n}=A_{m}B_{m} \]


Llegamos a: \vspace{0.5cm}

\begin{center}
\includegraphics[width=1\linewidth]{4.png} % <-- Sustituye esta ruta
\\[2mm]
\textbf{Figure 4}
\end{center}	

Tenemos: \vspace{0.5cm}

\begin{center}
\includegraphics[width=1\linewidth]{5.png} % <-- Sustituye esta ruta
\\[2mm]
\textbf{Figure 4}
\end{center}	


Con esto concluimos con el 



\end{document}